\documentclass[11pt,a4paper]{article}

\ifdefined\pdfinfoomitdate\pdfinfoomitdate=1\fi
\ifdefined\pdftrailerid\pdftrailerid{}\fi
\ifdefined\pdfsuppressptexinfo\pdfsuppressptexinfo=-1\fi

\usepackage[T1]{fontenc}
\usepackage[utf8]{inputenc}
\usepackage{lmodern}
\usepackage[a4paper,margin=28mm]{geometry}
\usepackage{amsmath}
\usepackage{booktabs}
\usepackage[colorlinks=true,linkcolor=blue]{hyperref}

\title{A Synthetic Note on Reproducible LaTeX Revision}
\author{LazyingArt project-owned sample}
\date{2 September 2026}

\begin{document}

\maketitle

\begin{abstract}
This synthetic document tests a small, reproducible LaTeX revision. It contains no
customer material and makes no scientific claim.
\end{abstract}

\section{Method}

A manuscript revision should freeze complete old and revised source trees and
build each one independently. The checked comparison state is represented by

\begin{equation}
  R = (S_{\mathrm{old}}, S_{\mathrm{new}}),
  \label{eq:build-state}
\end{equation}

where the two $S$ terms are named source snapshots compiled with the same engine
and dependency set.

\begin{table}[ht]
  \centering
  \caption{Artifacts checked in the revision workflow.}
  \label{tab:artifacts}
  \begin{tabular}{ll}
    \toprule
    Artifact & Check \\
    \midrule
    Baseline & Builds \\
    Revision & Builds \\
    \bottomrule
  \end{tabular}
\end{table}

\section{Conclusion}

A successful clean build is necessary but not sufficient. The two source
snapshots in \eqref{eq:build-state} remain separate and independently buildable.

\end{document}
